\originalchapter{哈密顿力学与正则变换}
\originalsection{Legendre变换}
拉格朗日方程把位置与速度联系起来; 哈密顿形式改用位置与共轭动量, 将二阶方程变为相空间内的一阶方程. 这个改写有条件: 速度Hessian$W_{ij}=L_{\dot q_i\dot q_j}$必须在所讨论区域可逆, 才能局部从$p_i=L_{\dot q_i}$解出$\dot q_i=v_i(q,p,t)$. 约束电磁场等奇异拉格朗日系统需其他理论, 本册不把它们硬套入正则变换.

定义$H(q,p,t)=\sum_ip_iv_i-L(q,v,t)$. 微分时$\sum_ip_i\dd v_i$与$\sum_iL_{v_i}\dd v_i$抵消, 故$\dd H=\sum_iv_i\dd p_i-\sum_iL_{q_i}\dd q_i-L_t\dd t$. 联合Euler--Lagrange方程得
\begin{equation}\label{eq:ham}
\dot q_i=H_{p_i},\qquad\dot p_i=-H_{q_i},\qquad H_t=-L_t.
\end{equation}
$H$是能量函数的相空间表达, 仅在通常固定约束及速度无关势时才自动等于$T+U$. 系统不显含时间时, 沿\eqref{eq:ham}有$\dot H=H_t=0$.

\originalsection{Poisson括号}
本册统一定义$\{f,g\}=\sum_i(f_{q_i}g_{p_i}-f_{p_i}g_{q_i})$, 所以$\{q_i,p_j\}=\delta_{ij}$, $\{p_i,q_j\}=-\delta_{ij}$. 两个相空间函数沿运动的演化为$\dot f=\{f,H\}+f_t$. 由乘积求导可逐项核对反对称性、线性和乘积法则. 展开二阶偏导并用混合偏导交换, 可核对Jacobi恒等式$\{f,\{g,h\}\}+\{g,\{h,f\}\}+\{h,\{f,g\}\}=0$. 因而两个不显含时间的守恒量$f,g$的括号也守恒: 由Jacobi式和$\{f,H\}=\{g,H\}=0$得$\{\{f,g\},H\}=0$.

将坐标排序为$z=(q_1,\ldots,q_n,p_1,\ldots,p_n)^{\mathsf T}$, 定义$\mathbb J=\left(\begin{smallmatrix}0&I\\-I&0\end{smallmatrix}\right)$. 则$\dot z=\mathbb J\nabla H$, $\{f,g\}=\nabla f^{\mathsf T}\mathbb J\nabla g$. 若新坐标$Z=Z(z)$的Jacobi矩阵为$A$, 保持括号等价于$A\mathbb JA^{\mathsf T}=\mathbb J$, 称为辛条件. 取行列式只先得$(\det A)^2=1$; 辛变换实际保定向, 可由辛体积形式$\omega^n/n!$不变得$\det A=1$. 在连续接近恒等的变换上也可直接选正支.

哈密顿流的散度为$\sum_i(H_{p_iq_i}-H_{q_ip_i})=0$. 体积元随流的Jacobi行列式满足$\dot D=(\operatorname{div}\dot z)D$, 因而$D=1$（初值为1）, 即Liouville定理. 它描述相体积, 并不意味着各坐标距离或各局部面积投影不变.

\originalsection{生成函数}
正则变换使作用量中的一形式之差为全微分:
$\sum_ip_i\dd q_i-H\dd t=\sum_iP_i\dd Q_i-K\dd t+\dd F$. 若以$(q,P,t)$为独立变量并令$F_2=F+\sum_iQ_iP_i$, 比较系数得
\[
p_i=\pd{F_2}{q_i},\qquad Q_i=\pd{F_2}{P_i},\qquad
K=H+\pd{F_2}{t}.
\]
$F_2=\sum_iq_iP_i$给恒等变换; 它说明生成函数的偏导分别承担新旧坐标的联系. 要真正定义局部可逆变换, 混合Hessian$F_{2,qP}$仍需可逆. 其他类型由Legendre变换得到:
\begin{align*}
F_1(q,Q):&\quad p=F_{1,q},\quad P=-F_{1,Q},\\
F_3(p,Q):&\quad q=-F_{3,p},\quad P=-F_{3,Q},\\
F_4(p,P):&\quad q=-F_{4,p},\quad Q=F_{4,P}.
\end{align*}
各型均有$K=H+F_{a,t}$, 其独立变量与可逆区域必须明确. 不能在$F_4$的偏导中把$q$也当独立变量.

\begin{example}
一维谐振子$H=p^2/(2m)+m\omega^2q^2/2$. 令$q=\sqrt{2I/(m\omega)}\sin\theta$, $p=\sqrt{2m\omega I}\cos\theta$. 求新方程.
\end{example}
\begin{solution}
直接求导得到$\{\theta,I\}=1$, 所以$(\theta,I)$是正则对. 也可算$p\dd q=I\dd\theta+\dd(I\sin\theta\cos\theta)$: 将$p\dd q$展开为$2I\cos^2\theta\dd\theta+\sin\theta\cos\theta\dd I$, 与右端一致. 新哈密顿量为$K=\omega I$, 故$\dot I=0$, $\dot\theta=\omega$. 原相轨道为椭圆, 在新坐标中是固定$I$上的匀速角运动. $I=0$处相位无定义, 因此此图不能覆盖相空间原点.
\end{solution}
\originalsection{速度线性项与质量矩阵}
固定坐标下常见$L=\dot q^{\mathsf T}M(q,t)\dot q/2+A(q,t)^{\mathsf T}\dot q-U(q,t)$, $M$对称可逆. 动量为$p=M\dot q+A$, 必须先解出$\dot q=M^{-1}(p-A)$, 再做Legendre变换, 得
\[
H=\frac12(p-A)^{\mathsf T}M^{-1}(p-A)+U.
\]
线性速度项不直接贡献能量函数中的二次动能, 却改变动量与速度的联系. 电磁场中$L=m\vect v^2/2+e\vect A\cdot\vect v-e\phi$, $\vect p=m\vect v+e\vect A$, $H=(\vect p-e\vect A)^2/(2m)+e\phi$. 由Euler--Lagrange式直接展开$\partial_iA_j-\partial_jA_i$, 得$m\dot{\vect v}=e(\vect E+\vect v\times\vect B)$, 其中$\vect E=-\nabla\phi-\partial_t\vect A$, $\vect B=\nabla\times\vect A$. 磁力不作功, 但共轭动量不等于机械动量.

规范变换$\vect A\mapsto\vect A+\nabla\chi$, $\phi\mapsto\phi-\partial_t\chi$使$L\mapsto L+e\dd\chi/\dd t$, 不改运动; 共轭动量变成$\vect p+e\nabla\chi$, 哈密顿表达还出现$-e\partial_t\chi$. 所以应比较同一机械状态, 不能用未变换的$p$数值比较两规范. 循环坐标对应的是所选规范中的共轭动量守恒, 若$L$仅差全导数而不严格不变, 应使用Noether端点项修正.

均匀磁场取$\vect A=\vect B\times\vect r/2$, 则$\nabla\times\vect A=\vect B$; 原8.223第14讲第9页文字的叉积次序反了, 其分量已用此正确次序. 一个具体例子是$\vect B=B\hat x$、重力$-mg\hat z$, 令$\beta=eB/m\ne0$. Lorentz方程给$\ddot y=\beta\dot z$, $\ddot z=-g-\beta\dot y$, $\ddot x=0$. 从原点静止出发, 先由$\ddot y=\beta\dot z$积分为$\dot y=\beta z$, 再解$\ddot z+\beta^2z=-g$, 得
\[
z=\frac g{\beta^2}(\cos\beta t-1),\qquad
y=\frac g{\beta^2}\sin\beta t-\frac g\beta t.
\]
运动叠加回旋和平均$y$方向漂移$-g/\beta$, 等于一般常力漂移$\vect v_D=\vect F\times\vect B/(eB^2)$. 初位置和速度、两条方程均可直接代回检验. 磁场不作功, 但重力作功, 所以不能误说速度大小恒定. 原稿第11页的积分常数与静止初值不一致, 这里据初值重算. 若粒子限制于理想固定导线, 磁力与切向速度垂直, 对沿线的广义坐标不作虚功; 这依赖理想约束反力, 不能推出无约束粒子轨道不变.

第1章的小车摆有
\[
M_q=\begin{pmatrix}M+m&m\ell\cos\theta\\m\ell\cos\theta&m\ell^2\end{pmatrix},
\quad \det M_q=m\ell^2(M+m\sin^2\theta)>0,
\]
所以对$M,m>0$处处可做Legendre变换. 其逆矩阵为$[m\ell^2(M+m\sin^2\theta)]^{-1}\left(\begin{smallmatrix}m\ell^2&-m\ell\cos\theta\\-m\ell\cos\theta&M+m\end{smallmatrix}\right)$, 故
\[
H=\frac{m\ell^2p_x^2+(M+m)p_\theta^2-2m\ell\cos\theta\,p_xp_\theta}
{2m\ell^2(M+m\sin^2\theta)}-mg\ell\cos\theta.
\]
由$\dot x=H_{p_x}$和$\dot\theta=H_{p_\theta}$重得速度矩阵逆变换, $\dot p_x=0$则重得水平动量守恒. 原8.09第11页的逆矩阵分母排印与此不一致, 以行列式及相乘得到单位阵核验. 不必将复杂$p_\theta$导数展开成难读长式, 可用矩阵表达核查与拉格朗日方程完全等价.

\originalsection{循环自由度的Routh约化}
消去循环坐标的速度时, 直接把守恒动量代回原$L$通常会错. 以平面中心力为例, $L=m(\dot r^2+r^2\dot\phi^2)/2-U(r)$, 固定$p_\phi=J$, 解$\dot\phi=J/(mr^2)$. 对这个速度先作部分Legendre变换$R=L-J\dot\phi$, 再代回, 得
\[
R(r,\dot r;J)=\frac m2\dot r^2-U(r)-\frac{J^2}{2mr^2}.
\]
变分$r$在固定$J$下进行, 给$m\ddot r=-U'(r)+J^2/(mr^3)$, 正确恢复离心项. 若直接代回$L$, 最后一项符号相反, 会错误地得到向内的离心项. 一般分成非循环$q$和循环$s$, 定义$R=L-p_s\dot s$, 在循环速度Hessian可逆时解$\dot s(q,\dot q,p_s,t)$, 对$q$得到正确方程. 不同书可能把$R$定义成相反符号, 须连同其运动方程检查.

\originalsection{相空间作用量与无穷小正则变换}
Hamilton方程也来自一阶作用量$S[q,p]=\int(p_i\dot q_i-H)\dd t$. 独立变分$q,p$, 对$p_i\delta\dot q_i$分部积分后
\[
\delta S=[p_i\delta q_i]_{t_1}^{t_2}
+\int\bigl[(\dot q_i-H_{p_i})\delta p_i
 +(-\dot p_i-H_{q_i})\delta q_i\bigr]\dd t.
\]
固定位置端点使边界项消失, 任意内部变分得到两条Hamilton方程. 这里没有要求同时固定所有动量端点, 否则会过度规定典型一阶初值问题.

令第二型生成函数$F_2=q_iP_i+\epsilon G(q,P)$, 到一阶有$Q_i=q_i+\epsilon G_{p_i}$、$P_i=p_i-\epsilon G_{q_i}$. 因而任意函数的变化为$\delta f=\epsilon\{f,G\}$. $G=p_j$生成第$j$方向平移; $G=\vect n\cdot(\vect q\times\vect p)$生成绕$\vect n$旋转, 给$\delta\vect q=\epsilon\vect n\times\vect q$、$\delta\vect p=\epsilon\vect n\times\vect p$. 分量相乘可核对$\{J_i,J_j\}=\epsilon_{ijk}J_k$. $G=H$生成沿Hamilton流的一小步时间演化. 若$G$不显含时间且$\{G,H\}=0$, 则这一个参数的变换保持$H$, 同时$G$守恒, 将连续对称与动力学联系起来.

点变换$Q_j=f_j(q,t)$也有生成函数$F_2=\sum_jf_j(q,t)P_j$, 给$p_i=\sum_jP_j\partial f_j/\partial q_i$与$K=H+\sum_jP_j\partial_tf_j$. 这里新旧动量用转置Jacobi矩阵联系, 指标不能错位; 原稿第64页的印刷式应按上述偏导核对.


本章来源: 8.223第15--18讲与8.09第1、4章. 本册的括号方向与8.223作业4所用方括号相反, 作业译写处会说明, 不混用约定. 规范、Routh约化及相空间变分的边界条件解释为中文补充.
